\chapter{Перенесення між машинами}\label{ch:g06}

Цей розділ~--- історія гіпотези, яку спростували з користю. Гіпотеза H5 припускала, що провал перенесення DIII-D$\to$MAST має топологічну складову. Вимір показав протилежне (R7), а звідси вийшли два встановлені факти: критерій машинно-незалежний (E15), і топологія переживає перенесення там, де значення не переживають (E16). Нижче також розглянуто колишню гіпотезу C8 (2026-09-29 спростовано як артефакт знака~--- R9), відкладені C9 і C10, ліцензійне питання Q13 і ліцензії E28.

\section{Звідки H5}\label{sec:g06-h5}

Наївне перенесення між машинами в пілоті челенджу падає з SSIM 0,83 до 0,10 (\texttt{01-landscape/\allowbreak saturated-vs-\allowbreak headroom.md}, позначка [V]~\cite{g06Headroom}). Зазвичай це читають так: «значення не переносяться». H5 ставила гостріше питання (docstring \texttt{h5\_transfer\_topology.py}):

\begin{quote}\itshape
чи гірша топологія міжмашинного передбачення за топологію внутрішньомашинного передбачення тієї самої $L^2$-точності?
\end{quote}

Якби так, модель переносила б \emph{значення}, але не \emph{скелет}, і топологічний член розд.~\ref{ch:g05} побачив би відмову, яку $R^2$ сам по собі недооцінює.

\paragraph{Модель замість моделі.} Лінійний декодер, навчений на DIII-D, може видати лише поле з лінійної оболонки PCA-базису DIII-D. Тож проєкція \emph{істинного} MAST на цей базис дає \emph{верхню межу} того, що будь-який такий декодер міг би вивести для MAST, причому за ідеального знання цілі. Якщо навіть ця проєкція втрачає топологію, жодна така модель її не збереже. Контроль~--- проєкція DIII-D на власний базис того самого розміру.

\section{Конвенції знака: \texttt{AXIS\_SIGN}}\label{sec:g06-sign}

\code{fec/starter/fusion_scoring/common.py}{74}{89}{(знакові конвенції скорера FEC)}

Рядки 74--80: на DIII-D магнітна вісь лежить у \emph{мінімумі} збереженого $\psi$ (99,98\,\% з 1\,559\,340 кадрів), а в EFIT++/FAIR-MAST~--- у \emph{максимумі} (100\,\% кадрів MAST). Тому \texttt{AXIS\_SIGN}$=\{-1,+1\}$, і скорер скрізь працює з добутком $\texttt{AXIS\_SIGN}\cdot\psi$, у якому вісь завжди є максимумом. Коментар підкреслює, що машини мають \emph{однаковий} напрям струму: різниця походить від двох реалізацій EFIT, а не від фізики. Рядки 82--89: геометрично ідеальне передбачення для MAST у конвенції DIII-D давало б $R^2_\psi=-5{,}71$, гірше за сталу. Тому з версії v3.1 скорер інваріантний до знака. Він обирає \emph{один біт на машину на весь фолд}, тож знаком не можна «фармити» шум.

У топології знак нічого не змінює: і максимум, і мінімум мають індекс $+1$ (розд.~\ref{ch:g05}). Тому критерій повинен працювати на обох машинах без жодної поправки. Саме це перевіряє перший блок скрипта.

\section{Де це в коді}\label{sec:g06-code}

\code{ourtry/04-novelty/h5_transfer_topology.py}{61}{84}{(\texttt{skeleton}: область і попіксельні індекси)}
\code{ourtry/04-novelty/h5_transfer_topology.py}{85}{97}{(продовження: компоненти, O/X, сума)}

Це самодостатня копія детектора з \texttt{topology\_probe.py} (розд.~\ref{ch:g05}), параметризована машиною. Рядок 64 викликає \texttt{extract\_lcfs} з назвою машини, тобто з її \texttt{AXIS\_SIGN}. Якщо LCFS не знайдено, функція повертає \texttt{None} (рядки 65--66): це «провал вилучення», який з'явиться в C8. Рядки 69--72 будують область з ерозією на дві комірки і відкидають області менші за 40 пікселів. Рядки 73--84 рахують попіксельне число обертів, рядки 85--97~--- зв'язні компоненти (виправлення V-D4) і лічильники $N_O$, $N_X$ та суму індексів.

\code{ourtry/04-novelty/h5_transfer_topology.py}{128}{138}{(\texttt{project}: найкращий вихід DIII-D-декодера)}

Рядок 134: $\mathrm{rec}=\mu+(Y_s-\mu)V_k^{\mathsf T}V_k$~--- ортогональна проєкція на перші $k$ головних компонент DIII-D навколо середнього поля DIII-D $\mu$. Рядки 132--137 для MAST пробують обидва глобальні знаки $s=\pm1$ і лишають той, що дає меншу $L^2$-похибку, як і скорер. Для DIII-D (контроль) знак не шукається.

\code{ourtry/04-novelty/h5_transfer_topology.py}{147}{167}{(головний цикл по розміру базису; лічильник \texttt{nfail})}

Для $k=5,10,20,50$ рядки 148--155 обчислюють $R^2_\psi$ і медіанний скелет контролю на 8 кадрах. Рядки 157--162 роблять те саме для MAST, де $R^2$ рахується відносно \emph{знакоскоригованої} цілі $s_M\psi_M$. Рядок 160 рахує \texttt{nfail}~--- кількість кадрів, на яких \texttt{skeleton} повернув \texttt{None}. Рядки 164--166 друкують рядок таблиці з приміткою «[$n$/8 LCFS extraction failed]».

\section{Що вийшло: E15, E16, R7}\label{sec:g06-results}

Повторний запуск 2026-09-29 (3 демо-розряди на машину, 500 кадрів DIII-D і 115 кадрів MAST) відтворив таблицю:

\begin{center}\small
\begin{tabular}{r|cc|ccc}
\toprule
$k$ & $R^2_\psi$ D3D$\to$D3D & скелет & $R^2_\psi$ MAST$\to$D3D & скелет & знак $s_M$\\
\midrule
5  & 0,99987 & 1O/0X & 0,66248 & 1O/0X [3/8 провалів LCFS] & $-1$\\
10 & 1,00000 & 1O/0X & 0,87883 & 1O/0X & $+1$\\
20 & 1,00000 & 1O/0X & 0,92690 & 1O/0X [2/8 провалів LCFS] & $-1$\\
50 & 1,00000 & 1O/0X & 0,96298 & 1O/0X & $+1$\\
\bottomrule
\end{tabular}
\end{center}

\noindent На істині обидві машини мають 1 еліптичну точку, 0 гіперболічних і суму $+1$ на 8/8 кадрах.

\begin{established}{E15, E16}
\textbf{E15.} Топологічний критерій машинно-незалежний: DIII-D і MAST дають 1O/0X і суму $+1$, попри протилежні конвенції знака, різні сітки й різну геометрію.
\textbf{E16.} Топологія переживає перенесення там, де значення не переживають: у проєкції MAST на базис DIII-D $R^2_\psi$ падає до 0,66, а скелет лишається 1O/0X~\cite{RegEst}. (Примітку про рядок $k=50$ зареєстровано 2026-09-29, див. нижче.)
\end{established}

\begin{refuted}{R7 (H5)}
«Провал перенесення має топологічну складову». Топологія вціліла: при $R^2_\psi=0{,}66$ скелет 1O/0X. Інверсія виявилася кориснішою за саму гіпотезу, бо дала E15 і E16~\cite{RegEst}. У підсумку реєстру гіпотез R7 названо одним із двох спростувань, що дали «кращі результати після спростування», ніж містили до нього~\cite{RegConj}.
\end{refuted}

\begin{derivation}[чому гладка похибка не народжує критичних точок (оцінка путівника)]
Нова критична точка з'являється там, де похибка компенсує градієнт: $\grad\delta\psi=-\grad\psi$. Проєкція на кілька головних компонент дає \emph{гладку} похибку з малими $|\grad\delta\psi|$. Скрізь, крім малого околу осі, $|\grad\psi|$ більший, тож нових нулів немає. Поблизу осі невироджений екстремум (функція Морса) структурно стійкий щодо $C^1$-малих збурень: він лише зсувається, але не розщеплюється. Отже, топологію захищає мала похибка \emph{в $C^1$}, а не в $L^2$. Білий шум на сітці з кроком $h$ має $|\grad\delta\psi|\sim\varepsilon/h$ і тому народжує пари O--X (розд.~\ref{ch:g05}). Проєкція ж міжмашинна, велика в $L^2$, але гладка. Саме тому E16 не суперечить E10.
\end{derivation}

\paragraph{Застереження, знайдене при підготовці путівника.} Скрипт друкує \emph{медіани} по 8 кадрах з форматом \texttt{.0f}, а Python округлює 0,5 до парного, тобто до 0. Покадровий прогін (\texttt{guide/figures/g06\_h5\_frames.py}) для $k=50$ (MAST$\to$D3D) дав такі скелети: $(1,0,{+1})$ на чотирьох кадрах, $(1,1,0)$ на трьох і $(2,1,{+1})$ на одному. Отже, медіана $N_X=0{,}5$, а на половині кадрів в області є X-точка, хоча в таблиці стоїть «1O/0X». Для $k=10$ і для всіх контрольних рядків DIII-D усі кадри мають $(1,0,{+1})$; для $k=5$ усі 5 вилучених кадрів мають $(1,0,{+1})$, а для $k=20$ (6/8 вилучено) один кадр має $(1,1,0)$, медіана 0 (покадровий прогін \texttt{Our try/04-novelty/h5\_frames\_check.py}, 2026-09-29). E16 (твердження про $k=5$) це не зачіпає. Проте рядок $k=50$ показує, що топологічні дефекти при перенесенні \emph{можливі}, і H5 варто перевірити ще раз при \emph{порівнянному} $R^2_\psi$ (див. практикум). Зареєстровано 2026-09-29 як примітку до E16 (і посилання в R7); скрипт~--- \texttt{Our try/04-novelty/h5\_frames\_check.py}~\cite{RegEst,RegLog}.

\section{C8: провали вилучення LCFS}\label{sec:g06-c8}

\begin{conjecture}{C8 --- \emph{закрито 2026-09-29: спростовано як артефакт, тепер R9}}
\emph{Нижче~--- формулювання, яке діяло до 2026-09-29.} Провали вилучення LCFS на 3/8 і 2/8 кадрах MAST при перенесенні (E16)~--- \emph{систематичний сигнал}, а не шум. \emph{Спростує (переписано 2026-09-29):} на вибірці MAST з EFIT-рівновагами з FAIR-MAST (CC BY-SA 4.0) або іншого джерела, значно більшій за 3 демо-розряди, провали розподілені по $k$ випадково. \emph{Вартість:} $\approx2$~год обчислень \emph{плюс} невідома вартість отримання даних і ліцензійне обтяження BY-SA (Q13, V-C4). У нинішньому формулюванні гіпотеза нездійсненна: \texttt{efit\_psirz} для MAST на HF публікується лише для 3 демо-розрядів (N3)~\cite{RegConj}.
\end{conjecture}

\paragraph{Спостереження для путівника: провали збігаються зі знаком.} У таблиці вище провали є \emph{лише} в рядках, де пошук знака обрав $s_M=-1$ ($k=5$ і $k=20$). Ми перевірили це прямо тими самими функціями \texttt{load}, \texttt{skeleton} і тим самим базисом (\texttt{guide/figures/g06\_c8\_sign\_check.py}). Для кожного $k$ і обох знаків реконструкцію повертали в конвенцію MAST множенням на $s_M$, а потім шукали LCFS. Результат: зі знаком $+1$ провалів 0 для всіх $k$. Зі знаком $-1$ «як є» маємо 3, 2, 2 і 8 провалів з 8 (для $k=5,10,20,50$). Після повернення знака провалів 0 скрізь.

Механізм описано в docstring \texttt{extract\_lcfs\_with\_sign} скорера. Без явного \texttt{axis\_sign} функція спершу пробує конвенцію машини, а протилежний знак бере як запасний варіант за правилом «перший успіх, а не найкраще припасування». На інвертованій карті такий шлях дає хибні межі або провали. \texttt{skeleton} передає \texttt{"MAST"} без \texttt{axis\_sign}, тоді як реконструкція, обрана зі знаком $-1$, інвертована.

Отже, «систематичний сигнал» C8 виявився артефактом обробки знака в самому скрипті, а не властивістю перенесення; для перевірки вистачило 3 демо-розрядів, нових даних не знадобилося. Зареєстровано 2026-09-29 як \textbf{R9} (C8~→ спростовано, артефакт); скрипт~--- \texttt{Our try/04-novelty/c8\_sign\_check.py}~\cite{RegEst,RegConj,RegLog}. Сам \texttt{skeleton} не змінювали: передавати \texttt{axis\_sign} лишається задачею практикуму.

\section{Відкладені гіпотези і ліцензії}\label{sec:g06-deferred}

\begin{itemcard}{C9}{відкладено свідомо}
Перенесення між машинами можна покращити фізично коваріантним поданням: коефіцієнтами функцій Гріна замість пікселів (пор.~з архітектурою~\cite{McClenaghan2024}). Відкладено, бо це задача \emph{живого} змагання Sophelio. Беремо її як мотивацію, а не як власну роботу~\cite{RegConj}.
\end{itemcard}

\begin{itemcard}{C10}{відкладено свідомо}
Контамінація міток: частина похибки будь-якої ML-моделі~--- це успадковане зміщення EFIT, а не похибка моделі. Для перевірки потрібна незалежна істина, якої немає. Записано як відкрите питання поля~\cite{RegConj}. Зв'язок із цим розділом: DIII-D і MAST реконструйовані \emph{різними} реалізаціями EFIT (\texttt{AXIS\_SIGN}), тож частина «міжмашинної» різниці може бути різницею реконструкцій.
\end{itemcard}

\begin{established}{E28 (прочитано)}
Ліцензії: HDB5~--- CC BY 4.0; ConStellaration~--- MIT; дані Sophelio~--- CC BY 4.0; FAIR-MAST~--- CC BY-SA 4.0, причому share-alike поширюється на похідні бази через §4(b)~\cite{RegEst,RegVerif}.
\end{established}

\begin{itemcard}{Q13}{відкрите, блокує перевипуск частини MAST}
Яка правова підстава переходу 1\,206 розрядів MAST із CC BY-SA 4.0 (FAIR-MAST) у випуск CC BY 4.0 (Sophelio)? Share-alike поширюється на похідні бази через §4(b), тож перевипуск як простий CC BY сама BY-SA не дозволяє. Адресат~--- Sophelio та/або UKAEA. До з'ясування частину MAST трактуємо як обтяжену BY-SA; частину DIII-D це не зачіпає~\cite{RegQ,RegVerif}. Переписана 2026-09-29 версія C8 потребувала саме даних FAIR-MAST і впиралася в це питання; після закриття C8 як R9 нові дані для неї не потрібні.
\end{itemcard}

\begin{practicum}[перенесення і топологія]
\textbf{Відтворити} (з \texttt{fusion equilibrium challenge/starter}, $\approx$1--2~хв): \texttt{.venv/bin/python "../../Our try/04-novelty/h5\_transfer\_topology.py"}.

\textbf{Задача~1 (C8 без нових даних).} У \texttt{skeleton} передайте в \texttt{extract\_lcfs} явний \texttt{axis\_sign=s\_M*AXIS\_SIGN["MAST"]} (сигнатуру доведеться розширити). Чи зникнуть провали? (Відповідь для перевірки: так, 0 провалів при всіх $k$~--- R9.)

\textbf{Задача~2 (H5 чесно).} Замість медіан друкуйте покадрові трійки $(N_O,N_X,\Sigma)$. Потім зрівняйте $R^2_\psi$: для контролю DIII-D додайте шум або зменште $k$ до 1--3, щоб отримати $R^2\approx0{,}96$, і порівняйте частку кадрів з внутрішньою X-точкою з рядком MAST при $k=50$. Чи з'являється топологічна складова H5 при порівнянній точності?

\textbf{Задача~3 (знак).} Поясніть, чому оптимальний $s_M$ стрибає між $-1$ і $+1$ зі зміною $k$ (різниця похибок при $k=10$: 114 проти 126). Чим це відрізняється від правила скорера «один біт на фолд»?
\end{practicum}
